\documentclass[11pt]{article}
\usepackage{amsmath,amssymb,bm}

\begin{document}
\section{Analytical model: disk embedded in infinte membrane}
\subsection{Tensionless membrane: analytical solution}

We consider the axisymmetric membrane in the absence of surface tension,
\(\gamma=0\).
The outer membrane energy reduces to
\begin{equation}
F_{\rm out}
=
\pi\int_R^\infty r\,dr
\Big[
\kappa (J-D\theta^d)^2
+\kappa (J+D\theta^p)^2
+\kappa_t \big[(\theta^d)^2+(\theta^p)^2\big]
\Big],
\end{equation}
where
\(
J=\phi'+\phi/r
\)
and
\(
D\theta=\theta'+\theta/r
\).

\paragraph{Euler--Lagrange identities.}
For \(\gamma=0\), variation with respect to \(\theta^d\) and \(\theta^p\) gives
\begin{equation}
\theta^p(r)=\theta^d(r)\equiv\theta(r),
\end{equation}
so the distal and proximal monolayers have identical tilt everywhere.
With this identification, the outer energy decouples into shape and tilt parts:
\begin{equation}
(J-D\theta)^2+(J+D\theta)^2
=
2J^2+2(D\theta)^2.
\end{equation}

\paragraph{Outer shape field.}
The shape enters only through \(J\).
Minimization of the outer energy yields
\begin{equation}
J(r)=\phi'(r)+\frac{\phi(r)}{r}=0,
\qquad r\ge R.
\end{equation}
The solution that decays at infinity is
\begin{equation}
\phi(r)=\phi_*\frac{R}{r},
\end{equation}
where \(\phi_*=\phi(R)\).
Integrating \(\phi=dz/dr\) gives a logarithmic profile,
\begin{equation}
z(r)-z(R)=\phi_*R\ln\!\frac{r}{R},
\end{equation}
so the tensionless membrane forms a slowly flaring ``trumpet'' rather than a cone.

\paragraph{Outer tilt field.}
The common tilt \(\theta(r)\) satisfies
\begin{equation}
r^2\theta''+r\theta'-(1+\lambda^2 r^2)\theta=0,
\qquad
\lambda=\sqrt{\frac{\kappa_t}{\kappa}}.
\end{equation}
Imposing decay at infinity gives
\begin{equation}
\theta(r)=\theta(R)\frac{K_1(\lambda r)}{K_1(\lambda R)},
\qquad r\ge R.
\end{equation}

\paragraph{Inner (disk-covered) region.}
Over the disk (\(0\le r\le R\)) the membrane is flat.
The tilt field \(\theta_C(r)\) minimizes
\begin{equation}
F_{\rm in}
=
2\pi\int_0^R r\,dr
\left[
\frac{\kappa}{2}(D\theta_C)^2
+\frac{\kappa_t}{2}\theta_C^2
\right],
\end{equation}
and satisfies the same equation as above.
Regularity at \(r=0\) gives
\begin{equation}
\theta_C(r)=\theta_B\frac{I_1(\lambda r)}{I_1(\lambda R)},
\end{equation}
where \(\theta_B=\theta_C(R)\).

\paragraph{Boundary matching (1D-style continuation).}
Director continuity at the rim implies
\begin{equation}
\theta^d(R)=\theta_B-\phi_*,
\qquad
\theta^p(R)=\phi_*.
\end{equation}
Since \(\theta^d=\theta^p\) in the tensionless minimum, this fixes
\begin{equation}
\phi_*=\frac{\theta_B}{2},
\qquad
\theta(R)=\frac{\theta_B}{2}.
\end{equation}

\paragraph{Total energy.}
The elastic energies are
\begin{align}
F_{\rm in,el}
&=
\pi\kappa R\lambda
\frac{I_0(\lambda R)}{I_1(\lambda R)}\theta_B^2,\\
F_{\rm out,el}
&=
\frac{\pi}{2}\kappa R\lambda
\frac{K_0(\lambda R)}{K_1(\lambda R)}\theta_B^2,
\end{align}
and the contact contribution is
\begin{equation}
F_{\rm cont}=-2\pi R\,h\frac{\Delta\varepsilon}{a}\,\theta_B.
\end{equation}
Thus
\begin{equation}
F_{\rm tot}(\theta_B)
=
\pi\kappa R\lambda
\left[
\frac{I_0(\lambda R)}{I_1(\lambda R)}
+\frac{1}{2}\frac{K_0(\lambda R)}{K_1(\lambda R)}
\right]\theta_B^2
-2\pi R\,h\frac{\Delta\varepsilon}{a}\,\theta_B.
\end{equation}

\paragraph{Optimal boundary tilt.}
Minimization yields
\begin{equation}
\boxed{
\theta_B^{\rm opt}
=
\frac{h(\Delta\varepsilon/a)}
{\sqrt{\kappa\kappa_t}
\left[
\dfrac{I_0(\lambda R)}{I_1(\lambda R)}
+\dfrac{1}{2}\dfrac{K_0(\lambda R)}{K_1(\lambda R)}
\right]}
}
\end{equation}
and therefore
\begin{equation}
\phi_*^{\rm opt}=\frac{\theta_B^{\rm opt}}{2}.
\end{equation}

\paragraph{Optimal fields.}
The minimizing fields are
\begin{align}
\theta_C^{\rm opt}(r)
&=\theta_B^{\rm opt}\frac{I_1(\lambda r)}{I_1(\lambda R)},
\qquad 0\le r\le R,\\
\theta_d^{\rm opt}(r)
&=\theta_p^{\rm opt}(r)
=\frac{\theta_B^{\rm opt}}{2}
\frac{K_1(\lambda r)}{K_1(\lambda R)},
\qquad r\ge R,\\
\phi^{\rm opt}(r)
&=\frac{\theta_B^{\rm opt}}{2}\frac{R}{r},
\qquad r\ge R.
\end{align}

\subsection{Finite tension: one-disk solution (axisymmetric, free rim)}

We consider a single disk of radius $R$ covered by a flat membrane patch ($0\le r\le R$) and
a free outer membrane ($r\ge R$) under finite surface tension $\gamma>0$.
We work in the small-angle regime, so the midplane slope is described by the tangent angle $\phi(r)$.
For convenience define
\begin{equation}
J(r)\equiv \phi'(r)+\frac{\phi(r)}{r},\qquad
D\theta(r)\equiv \theta'(r)+\frac{\theta(r)}{r},
\qquad
\lambda \equiv \sqrt{\frac{\kappa_t}{\kappa}}.
\end{equation}

%------------------------------------------------------------
\paragraph{Outer free membrane energy ($r\ge R$).}
The outer energy reads
\begin{equation}
F_{\rm out}
=\pi\int_R^\infty r\,dr
\Big[
\kappa\big(J-D\theta^d\big)^2+\kappa_t(\theta^d)^2
+\kappa\big(J+D\theta^p\big)^2+\kappa_t(\theta^p)^2
+\gamma \phi^2
\Big].
\end{equation}

Variation yields a coupled Euler--Lagrange system which implies the exact identity
\begin{equation}
\boxed{\ \theta^p(r)-\theta^d(r)=\frac{\gamma}{\kappa_t}\,\phi(r)\ }.
\label{eq:theta_diff_identity}
\end{equation}
Eliminating $\theta^{d,p}$ gives a closed equation for $\phi$ of the form
\begin{equation}
r^2\phi''+r\phi'-(1+\psi^2 r^2)\phi=0,
\qquad
\boxed{\ \psi^2=\frac{\gamma}{\kappa}\,\frac{\kappa_t}{2\kappa_t+\gamma}\ }.
\end{equation}
The decaying solution is
\begin{equation}
\boxed{
\phi(r)=\phi_*\frac{K_1(\psi r)}{K_1(\psi R)}
},
\qquad \phi_*\equiv \phi(R).
\label{eq:phi_outer}
\end{equation}
Moreover,
\begin{equation}
J(R)=\phi'(R)+\frac{\phi_*}{R}
=-\psi\,\phi_*\,\frac{K_0(\psi R)}{K_1(\psi R)}.
\end{equation}

The tilt fields solve forced modified-Bessel equations.
A convenient representation is
\begin{align}
\theta^d(r)&=A\,K_1(\lambda r)+q\,\phi(r),\\
\theta^p(r)&=A\,K_1(\lambda r)-q\,\phi(r),
\end{align}
where the particular-solution coefficient is
\begin{equation}
q=\frac{\psi^2}{\psi^2-\lambda^2}=-\frac{\gamma}{2\kappa_t},
\end{equation}
so that Eq.~\eqref{eq:theta_diff_identity} is satisfied identically.

%------------------------------------------------------------
\paragraph{Inner disk-covered region ($0\le r\le R$).}
Over the disk the membrane is flat, and the tilt field $\theta_C(r)$ minimizes
\begin{equation}
F_{\rm in}
=2\pi\int_0^R r\,dr
\left[
\frac{\kappa}{2}(D\theta_C)^2+\frac{\kappa_t}{2}\theta_C^2
\right],
\end{equation}
giving
\begin{equation}
\boxed{
\theta_C(r)=\theta_B\frac{I_1(\lambda r)}{I_1(\lambda R)}
},
\qquad \theta_B\equiv \theta_C(R).
\label{eq:thetaC_inner}
\end{equation}
The corresponding exact inner elastic energy is
\begin{equation}
\boxed{
F_{\rm in,el}
=\pi\kappa R\lambda\,\theta_B^2\,\frac{I_0(\lambda R)}{I_1(\lambda R)}
}.
\label{eq:Fin_el}
\end{equation}

%------------------------------------------------------------
\paragraph{Contact driving term.}
We include the linear ``contact'' contribution
\begin{equation}
F_{\rm cont}=-2\pi R\,h\,\frac{\Delta\varepsilon}{a}\,\theta_B.
\label{eq:Fcont}
\end{equation}

%------------------------------------------------------------
\paragraph{Free-rim kinematics (1D-style director continuation).}
To generalize the 1D building-block construction, we impose at $r=R$
\begin{equation}
\boxed{
\theta^d(R)=\theta_B-\phi_*,
\qquad
\theta^p(R)=\phi_*
}.
\label{eq:BC_kin}
\end{equation}
Using $\theta^{d,p}(R)=A K_1(\lambda R)\pm q\phi_*$ and $q=-\gamma/(2\kappa_t)$,
Eq.~\eqref{eq:BC_kin} yields
\begin{equation}
\boxed{
A=\frac{\theta_B}{2K_1(\lambda R)},
\qquad
\phi_*=\frac{\theta_B}{2\mu},
\qquad
\mu\equiv 1-\frac{\gamma}{2\kappa_t}.
}
\label{eq:A_phi_star_vs_thetaB}
\end{equation}

%------------------------------------------------------------
\paragraph{Outer energy reduced to a quadratic in $\theta_B$.}
Evaluating the outer energy on-shell yields an effective quadratic form in the boundary values.
Define the outer ``Dirichlet-to-Neumann'' stiffnesses
\begin{equation}
A_\psi \equiv \pi\kappa R\,\psi\,\frac{K_0(\psi R)}{K_1(\psi R)},
\qquad
B_\lambda \equiv \pi\kappa R\,\lambda\,\frac{K_0(\lambda R)}{K_1(\lambda R)}.
\end{equation}
Using Eq.~\eqref{eq:A_phi_star_vs_thetaB} one finds
\begin{equation}
\boxed{
F_{\rm out}
=\left[
\frac{A_\psi}{4\mu^2}+\frac{B_\lambda}{2}
\right]\theta_B^2
\;+\;\text{const}.
}
\label{eq:Fout_reduced_thetaB}
\end{equation}
Together with Eqs.~\eqref{eq:Fin_el} and \eqref{eq:Fcont}, the total energy becomes
\begin{equation}
F_{\rm tot}(\theta_B)
=
\underbrace{\left[
\pi\kappa R\lambda\frac{I_0(\lambda R)}{I_1(\lambda R)}
+\frac{1}{2}\pi\kappa R\lambda\frac{K_0(\lambda R)}{K_1(\lambda R)}
+\frac{1}{4\mu^2}\pi\kappa R\psi\frac{K_0(\psi R)}{K_1(\psi R)}
\right]}_{\equiv\,K_{\rm eff}}
\theta_B^2
-2\pi R\,h\,\frac{\Delta\varepsilon}{a}\,\theta_B
+\text{const}.
\end{equation}

\paragraph{Optimal boundary tilt and kink.}
Minimizing $F_{\rm tot}$ gives
\begin{equation}
\boxed{
\theta_B^{\rm opt}
=
\frac{h(\Delta\varepsilon/a)}{
\kappa\left[
\lambda\frac{I_0(\lambda R)}{I_1(\lambda R)}
+\frac{\lambda}{2}\frac{K_0(\lambda R)}{K_1(\lambda R)}
+\frac{\psi}{4\mu^2}\frac{K_0(\psi R)}{K_1(\psi R)}
\right]
}
},
\qquad
\boxed{
\phi_*^{\rm opt}=\frac{\theta_B^{\rm opt}}{2\mu}
}.
\label{eq:thetaBopt_phiopt}
\end{equation}

%------------------------------------------------------------
\paragraph{Optimal fields.}
The minimizing inner field is
\begin{equation}
\boxed{
\theta_C^{\rm opt}(r)
=
\theta_B^{\rm opt}\,\frac{I_1(\lambda r)}{I_1(\lambda R)},
\qquad 0\le r\le R.
}
\end{equation}
The minimizing outer shape field is
\begin{equation}
\boxed{
\phi^{\rm opt}(r)=\phi_*^{\rm opt}\,\frac{K_1(\psi r)}{K_1(\psi R)},
\qquad r\ge R.
}
\end{equation}
Finally, the minimizing distal/proximal tilts are
\begin{equation}
\boxed{
\theta_d^{\rm opt}(r)
=
\frac{\theta_B^{\rm opt}}{2K_1(\lambda R)}\,K_1(\lambda r)
-\frac{\gamma}{2\kappa_t}\,\phi^{\rm opt}(r),
\qquad
\theta_p^{\rm opt}(r)
=
\frac{\theta_B^{\rm opt}}{2K_1(\lambda R)}\,K_1(\lambda r)
+\frac{\gamma}{2\kappa_t}\,\phi^{\rm opt}(r)
}
\end{equation}
for the case of $r\ge R$.


These satisfy $\theta_p^{\rm opt}-\theta_d^{\rm opt}=(\gamma/\kappa_t)\phi^{\rm opt}$ as required.

\section{Numerical parameter choice and reference energy scale}
\label{sec:one_disk_numbers}

In order to provide a concrete benchmark for comparison with numerical
minimization, we now fix a specific set of parameters and evaluate the
optimized energy of a single caveolin disk embedded in an infinite,
tensionless membrane.

\subsection{Choice of units and parameters}

We choose the bending modulus as the energy unit,
\begin{equation}
\kappa \equiv 1 ,
\end{equation}
corresponding physically to a monolayer bending rigidity
\(\kappa = 10\,k_B T\).

Lengths are measured in simulation units such that
\begin{equation}
R_{\mathrm{sim}} = 0.466666\ldots \equiv \frac{7\,\mathrm{nm}}{15\,\mathrm{nm}},
\end{equation}
i.e. one simulation length unit corresponds to \(15\,\mathrm{nm}\).

The remaining parameters are chosen as
\begin{align}
\kappa_t &= 225 , \\
\lambda &= \sqrt{\kappa_t / \kappa} = 15 , \\
\frac{h \Delta \varepsilon}{a} &= 4.286 ,
\end{align}
which correspond physically to
\(\kappa_t = 10\,k_B T/\mathrm{nm}^2\),
\(h = 2\,\mathrm{nm}\),
\(a = 0.7\,\mathrm{nm}^2\),
and \(\Delta\varepsilon = 1\,k_B T\).

With these values, the key dimensionless combination is
\begin{equation}
\lambda R = 15 \times 0.466666\ldots = 7 .
\end{equation}

\subsection{Unscaled physical parametrization}

We now rewrite the one-disk benchmark directly in physical units, without the
intermediate rescaling by $15\,\mathrm{nm}$. We take
\begin{equation}
R = 7\,\mathrm{nm}, \qquad
h = 2\,\mathrm{nm}, \qquad
a = 0.7\,\mathrm{nm}^2, \qquad
\Delta\varepsilon = 1\,k_B T,
\label{eq:phys_params}
\end{equation}
and use the monolayer bending rigidity as the energy unit,
\begin{equation}
\kappa \equiv 1,
\qquad \text{with physical value} \qquad
\kappa = 10\,k_B T.
\label{eq:kappa_unit_unscaled}
\end{equation}
The tilt modulus is taken as
\begin{equation}
\kappa_t = 10\,k_B T/\mathrm{nm}^2.
\label{eq:kappat_unscaled}
\end{equation}

With this choice,
\begin{equation}
\lambda \equiv \sqrt{\frac{\kappa_t}{\kappa}}
= \sqrt{\frac{10\,k_B T/\mathrm{nm}^2}{10\,k_B T}}
= 1\,\mathrm{nm}^{-1},
\label{eq:lambda_unscaled}
\end{equation}
so that the key dimensionless combination is simply
\begin{equation}
\lambda R = 7.
\label{eq:lambdaR_unscaled}
\end{equation}
Moreover,
\begin{equation}
\frac{h\Delta\varepsilon}{a\kappa}
=
\frac{(2\,\mathrm{nm})(1\,k_B T)}
{(0.7\,\mathrm{nm}^2)(10\,k_B T)}
=
0.285714\,\mathrm{nm}^{-1}.
\label{eq:contact_ratio_unscaled}
\end{equation}

For a tensionless membrane, the total energy has the quadratic form
\begin{equation}
F_{\mathrm{tot}}(\theta_B)
=
A\,\theta_B^2 - B\,\theta_B,
\label{eq:Ftot_quad_unscaled}
\end{equation}
with
\begin{equation}
A
=
\pi \kappa R \lambda
\left[
\frac{I_0(\lambda R)}{I_1(\lambda R)}
+
\frac{1}{2}\frac{K_0(\lambda R)}{K_1(\lambda R)}
\right],
\label{eq:A_unscaled}
\end{equation}
and
\begin{equation}
B
=
2\pi R \left(\frac{h\Delta\varepsilon}{a}\right)\frac{1}{\kappa}
=
2\pi R \left(\frac{h\Delta\varepsilon}{a\kappa}\right).
\label{eq:B_unscaled}
\end{equation}

Substituting $\lambda R = 7$ gives
\begin{equation}
A
=
\pi \cdot 1 \cdot 7 \cdot 1
\left[
\frac{I_0(7)}{I_1(7)}
+
\frac{1}{2}\frac{K_0(7)}{K_1(7)}
\right],
\label{eq:A_eval_unscaled}
\end{equation}
\begin{equation}
B
=
2\pi \cdot 7 \cdot 0.285714.
\label{eq:B_eval_unscaled}
\end{equation}
Using
\begin{equation}
\frac{I_0(7)}{I_1(7)} \simeq 1.08046,
\qquad
\frac{K_0(7)}{K_1(7)} \simeq 0.93530,
\label{eq:bessel_ratios_unscaled}
\end{equation}
we obtain
\begin{equation}
A \simeq 34.04,
\qquad
B \simeq 12.57.
\label{eq:AB_values_unscaled}
\end{equation}

Minimization of Eq.~\eqref{eq:Ftot_quad_unscaled} gives the optimal boundary tilt
\begin{equation}
\theta_B^*
=
\frac{B}{2A}
\simeq 0.1846 \approx 0.185.
\label{eq:thetaBstar_unscaled}
\end{equation}

The corresponding inner elastic energy is
\begin{equation}
F_{\mathrm{in,el}}
=
\pi \kappa R \lambda
\frac{I_0(\lambda R)}{I_1(\lambda R)}
(\theta_B^*)^2,
\label{eq:Fin_unscaled}
\end{equation}
which evaluates to
\begin{equation}
F_{\mathrm{in,el}} \simeq 0.81.
\label{eq:Fin_value_unscaled}
\end{equation}

The outer elastic energy is
\begin{equation}
F_{\mathrm{out,el}}
=
\frac{\pi}{2}\kappa R \lambda
\frac{K_0(\lambda R)}{K_1(\lambda R)}
(\theta_B^*)^2,
\label{eq:Fout_unscaled}
\end{equation}
which gives
\begin{equation}
F_{\mathrm{out,el}} \simeq 0.35.
\label{eq:Fout_value_unscaled}
\end{equation}

The contact contribution is
\begin{equation}
F_{\mathrm{cont}}
=
-2\pi R \left(\frac{h\Delta\varepsilon}{a}\right)\frac{\theta_B^*}{\kappa}
=
-B\,\theta_B^*,
\label{eq:Fcont_unscaled}
\end{equation}
hence
\begin{equation}
F_{\mathrm{cont}} \simeq -2.32.
\label{eq:Fcont_value_unscaled}
\end{equation}

Finally, the optimized total energy is
\begin{equation}
F_{\mathrm{tot}}^*
=
F_{\mathrm{in,el}} + F_{\mathrm{out,el}} + F_{\mathrm{cont}}
\simeq -1.16,
\label{eq:Ftot_value_unscaled}
\end{equation}
in units of $\kappa$.

Thus, working directly with the physical radius $R=7\,\mathrm{nm}$ and
$\lambda = 1\,\mathrm{nm}^{-1}$ gives exactly the same dimensionless benchmark
as the rescaled formulation, because the controlling combination $\lambda R$
remains equal to $7$.

\subsection{Optimized boundary tilt}

For a tensionless membrane, the total energy reduces to a quadratic form
in the boundary tilt angle \(\theta_B\),
\begin{equation}
F_{\mathrm{tot}}(\theta_B)
=
A\,\theta_B^2 - B\,\theta_B ,
\end{equation}
with
\begin{align}
A &= \pi \kappa R \lambda
\left[
\frac{I_0(\lambda R)}{I_1(\lambda R)}
+
\frac{1}{2}
\frac{K_0(\lambda R)}{K_1(\lambda R)}
\right], \\
B &= 2\pi R \left( \frac{h \Delta\varepsilon}{a} \right).
\end{align}

Evaluating the Bessel-function ratios at \(\lambda R = 7\),
\begin{equation}
\frac{I_0(7)}{I_1(7)} \simeq 1.08046,
\qquad
\frac{K_0(7)}{K_1(7)} \simeq 0.93530 ,
\end{equation}
gives
\begin{equation}
A \simeq 34.04,
\qquad
B \simeq 12.57 .
\end{equation}

Minimization with respect to \(\theta_B\) yields
\begin{equation}
\theta_B^{\ast} = \frac{B}{2A} \simeq 0.185 .
\end{equation}

\subsection{Energy breakdown and total energy}

At the optimal tilt angle, the individual contributions to the energy are
\begin{align}
F_{\mathrm{in,el}} &\simeq +0.81 , \\
F_{\mathrm{out,el}} &\simeq +0.35 , \\
F_{\mathrm{contact}} &\simeq -2.32 .
\end{align}

Summing all contributions gives the optimized total energy
\begin{equation}
F_{\mathrm{tot}}^{\ast}
=
F_{\mathrm{in,el}} + F_{\mathrm{out,el}} + F_{\mathrm{contact}}
\simeq -1.16 ,
\end{equation}
in units of the bending modulus \(\kappa\).

This value provides a reference benchmark for numerical simulations of
a single caveolin disk embedded in an effectively infinite,
tensionless membrane.


\end{document}
